\unit{Bài trình bày 182 — Hình học hình thức và hình học đại số}
\sid{182}
\markboth{Thuyết trình 182}{Thuyết trình 182}
\src{193}
\seminar{(Tháng 5 năm 1959)}
\paperhead{HÌNH HỌC HÌNH THỨC VÀ HÌNH HỌC ĐẠI SỐ}%
  {do Alexander GROTHENDIECK}

\fsection{1}{Lược đồ}

Ta biết rằng một không gian đại số affine xác định trên một trường $k$ về căn
bản được xác định bởi đại số affine $A$ của nó (vành các hàm chính quy xác định
trên $k$), vì các cấu xạ giữa các không gian đại số
$X\longrightarrow Y$ tương ứng song ánh với các đồng cấu
$k$-đại số $A(Y)\longrightarrow A(X)$. Đại số affine tương ứng với một không
gian đại số là một $k$-đại số kiểu hữu hạn và, theo quan điểm
\fquote{cổ điển}, không có phần tử lũy linh; ngược lại, mọi đại số thuộc kiểu
này đều nhận được như đại số affine của một không gian đại số xác định trên
$k$. Như vậy có một từ điển quen thuộc cho phép diễn giải các tình huống liên
quan đến không gian đại số affine bằng ngôn ngữ đại số giao hoán. Từ lâu người
ta đã nhận thấy rằng cách làm này cho những mệnh đề tổng quát hơn, vì thông
thường không còn cần giả sử các vành đang xét thuộc kiểu vừa nêu; giả thiết
Noether thường là đủ. Đặc biệt, dù có cho trước một trường cơ sở hay không,
không có lý do gì để loại trừ trường hợp các vành ấy chứa phần tử lũy linh.
Cho đến nay, các nhà hình học vẫn từ chối lưu ý đến những chỉ dẫn này và khăng
khăng chỉ xét các đại số affine không có phần tử lũy linh, nghĩa là các không
gian đại số mà bó cấu trúc không có phần tử lũy linh (và phần lớn còn chỉ xét
các không gian đại số \fquote{bất khả quy tuyệt đối}). Báo cáo viên cho rằng
quan niệm ấy đã gây trở ngại nghiêm trọng cho sự phát triển của các phương
pháp thật sự tự nhiên trong hình học đại số.

Cho $A$ là một vành giao hoán. Ta biết rằng tập
$X=\operatorname{Spec}(A)$ các iđêan nguyên tố của $A$ mang một tôpô tự nhiên,
gọi là \fquote{tôpô Zariski} hay tôpô phổ. Mặt khác, trên $X$ có một bó các
vành giao hoán $\Osh_X$, với thớ tại $\mathfrak p\in X$ là vành địa phương hóa
$A_{\mathfrak p}$ và vành các tiết diện của nó được đồng nhất với $A$. Do đó
$X$ trở thành một \emph{không gian có vành}, gọi là \emph{phổ nguyên tố} của
$A$. Một đồng cấu vành $f:A\longrightarrow B$ xác định một cấu xạ của các
không gian có vành
$f':\operatorname{Spec}(B)\longrightarrow\operatorname{Spec}(A)$; ánh xạ tập
hợp nền chính là $\mathfrak p\longmapsto f^{-1}(\mathfrak p)$. Các cấu xạ

\src{194}
của các không gian có vành từ $\operatorname{Spec}(B)$ đến
$\operatorname{Spec}(A)$ thu được theo cách này chính là những cấu xạ mà các
đồng cấu $\Osh_x\longrightarrow\Osh_y$ ($x=f'(y)$) đều là đồng cấu địa phương
(nghĩa là nghịch ảnh của iđêan cực đại là iđêan cực đại).

Ta gọi một không gian có vành đẳng cấu với một
$\operatorname{Spec}(A)$ nào đó là \emph{lược đồ affine}, và gọi một không gian
có vành affine địa phương là \emph{tiền lược đồ}; nói cách khác, mỗi điểm của
nó có một lân cận mở mà, với cấu trúc cảm sinh, là một lược đồ affine. Các
\emph{cấu xạ} của tiền lược đồ được định nghĩa theo cách hiển nhiên; ở địa
phương chúng tương ứng với các đồng cấu vành.

Khi cố định một tiền lược đồ $S$ và xét các cấu xạ tiền lược đồ
$X\longrightarrow S$, không gian $S$ đóng vai trò của một trường hay một vành
cơ sở (hay đúng hơn, của một không gian cơ sở trong một phân thớ). Khi đó ta
nói $X$ là một $S$-tiền lược đồ; nếu $S=\operatorname{Spec}(A)$, điều này cũng
có nghĩa $\Osh_X$ là một bó các $A$-đại số. Vì thế, mỗi tiền lược đồ có thể
được xem một cách duy nhất như một $\mathbf Z$-tiền lược đồ. Dĩ nhiên, các
$S$-tiền lược đồ lập thành một phạm trù; hơn nữa, người ta chứng minh rằng
trong phạm trù này tích của hai vật $X,Y$ luôn tồn tại và được ký hiệu
$X\times_S Y$. Khái niệm tích này cho phép định nghĩa phép thay đổi cơ sở của
một $S$-tiền lược đồ ứng với một cấu xạ $S'\longrightarrow S$: quả vậy,
$X\times_S S'$ có thể được xem như một $S'$-tiền lược đồ.

Ta nói $X$ \emph{tách} trên $S$ nếu đường chéo của $X$ trong
$X\times_S X$ là đóng. Một tiền lược đồ tách trên $\mathbf Z$ được gọi là một
\emph{lược đồ}; khi ấy nó tách trên mọi cơ sở. Để đơn giản, từ đây ta chỉ nói
đến các lược đồ và còn giả sử chúng là \emph{Noether}, nghĩa là hợp hữu hạn
của các tập mở affine vốn là phổ của các vành Noether. Ta nói $X$ \emph{thuộc
kiểu hữu hạn trên} $S$ nếu với mọi tập mở affine $U$ của $S$, nghịch ảnh của
$U$ trong $X$ là hợp hữu hạn của các tập mở affine có các vành là đại số kiểu
hữu hạn trên vành của $U$. Chính những $S$-lược đồ như vậy mới thích hợp cho
một nghiên cứu thực sự mang tính hình học. Đặc biệt, với mọi $s\in S$, thớ
$f^{-1}(s)$ của $X$ trên $s$ là một lược đồ đại số trên trường thặng dư
$\kappa(s)$ của vành địa phương $\Osh_s$ tại $s$ trong $S$. Bởi vậy, trong một
chừng mực nào đó, có thể xem $X$ như một họ các \fquote{không gian đại số}
$f^{-1}(s)$, trong đó tham số $s$ chạy trên $S$ (nghĩa là, theo quan điểm địa
phương, trên tập các iđêan nguyên tố của một vành đã cho). Dĩ nhiên, các
$\kappa(s)$ có thể có những đặc số khác nhau. Nếu
$S=\operatorname{Spec}(k)$, với $k$ là một trường, ta về căn bản thu lại khái
niệm thông thường về \fquote{không gian đại số}, chỉ khác ở chỗ giờ đây bó cấu
trúc có thể có các phần tử lũy linh.

\src{195}
Lấy cảm hứng từ các khái niệm quen thuộc, ta định nghĩa khái niệm
\emph{cấu xạ xạ ảnh} và, tổng quát hơn, \emph{cấu xạ proper}. Một cấu xạ như
vậy thuộc kiểu hữu hạn; hơn nữa, nó biến các tập con đóng thành các tập con
đóng và giữ nguyên tính chất này sau mọi phép thay đổi cơ sở.

Với $X$ là một lược đồ (Noether, như vẫn giả sử), bó $\Osh_X$ là một
\emph{bó vành coherent} theo nghĩa của [2]. Do đó, các bó môđun coherent trên
$X$ cũng chính là những bó đẳng cấu địa phương với đối hạch của một đồng cấu
$\Osh_X^{m}\longrightarrow\Osh_X^{n}$.

\fsection{2}{Lược đồ hình thức}

Cho $X$ là một lược đồ và $X'$ là một tập con đóng của $X$. Khi đó tồn tại một
bó con coherent $J$ của $\Osh_X$ sao cho
$X'=\operatorname{Supp}(\Osh_X/J)$ (thậm chí còn có một bó lớn nhất như vậy).
Khi được trang bị $\Osh_X/J$, không gian $X'$ trở thành một lược đồ, ký hiệu
$X_0$; lược đồ ấy được gọi là \fquote{lược đồ con đóng của $X$}. Với mọi $n$,
ta cũng có thể xét $X'$ được trang bị $\Osh_X/J^{n+1}$, ký hiệu $X_n$; đây là
một tiền lược đồ con đóng của $X$, có tập nền vẫn là $X'$ nhưng mang một bó
cấu trúc khác, cụ thể là $\Osh_{X_n}=\Osh_X/J^{n+1}$. Hiển nhiên, các
$\Osh_{X_n}$ lập thành một hệ xạ ảnh các bó vành trên $X'$; giới hạn xạ ảnh
$\overline{\Osh}_X$ của hệ ấy được gọi là \emph{hoàn thành hình thức của
$\Osh_X$ dọc theo $X'$}. Khi được trang bị bó vành này, $X'$ được gọi là
\emph{hoàn thành hình thức của $X$ dọc theo $X'$}; đó là một không gian có
vành, nhưng nói chung không phải là một lược đồ. Với mọi bó coherent $F$ trên
$X$, tương tự ta có thể xét hoàn thành hình thức
$\overline F=\varprojlim_n F_n$ của $F$ dọc theo $X'$, trong đó
$F_n=F\otimes_{\Osh_X}\Osh_X/J^{n+1}$; đây là một bó môđun trên
$\overline X$. Các tiết diện của nó được gọi là \fquote{các tiết diện hình
thức của $F$ dọc theo $X'$} và được đồng nhất với các phần tử của
$\varprojlim_n\Gamma(X',F_n)$. Với $F=\Osh_X$, ta thu được các
\fquote{hàm chỉnh hình} của $X$ dọc theo $X'$ theo nghĩa Zariski (chúng tôi
không dùng thuật ngữ ấy vì nó xung đột với thuật ngữ cổ điển).

Ta gọi một không gian tôpô $\mathfrak X$, được trang bị một bó các vành tôpô
$\Osh_{\mathfrak X}$ thỏa mãn điều kiện sau, là một \emph{lược đồ hình thức}
(ngầm hiểu: Noether): có một đẳng cấu các bó vành tôpô
$\Osh_{\mathfrak X}=\varprojlim_n\Osh_n$, trong đó các $\Osh_n$ lập thành một
hệ xạ ảnh các bó vành trên $\mathfrak X$, mỗi bó biến $\mathfrak X$ thành một
lược đồ $\mathfrak X_n$, và với $m\geq n$, đồng cấu
$\Osh_m\longrightarrow\Osh_n$ là toàn ánh, có hạch bằng $J_m^{n+1}$, trong đó
$J_m$ là

\src{196}
hạch của $\Osh_m\longrightarrow\Osh_0$. Khi đó người ta chứng minh được rằng
$\Osh_{\mathfrak X}$ là một \emph{bó coherent các vành địa phương Noether}.

Theo các định nghĩa, một hoàn thành hình thức $\overline X$ như trên là một
lược đồ hình thức; ngược lại, mọi lược đồ hình thức đều \emph{ở địa phương}
thuộc dạng này. Thật vậy, cho một lược đồ hình thức \emph{affine} (nghĩa là
$\mathfrak X_0$ affine, kéo theo mọi $\mathfrak X_n$ đều affine) tương đương
với cho một vành tôpô Noether, tách và đầy đủ đối với tôpô $J$-adic.

Các định nghĩa thông thường về cấu xạ, cấu xạ kiểu hữu hạn, cấu xạ proper,
v.v.\ của các lược đồ thông thường được mở rộng không khó khăn sang các lược
đồ hình thức.

\fsection{3}{Ba định lý cơ bản}

Cho $f:X\longrightarrow Y$ là một cấu xạ proper của các lược đồ (Noether, như
vẫn giả sử), $Y'$ là một tập con đóng của $Y$, còn $X'$ là nghịch ảnh của nó
trong $X$; xét các hoàn thành hình thức tương ứng $\overline Y$ và
$\overline X$. Khi đó $f$ cảm sinh một cấu xạ
$\overline f:\overline X\longrightarrow\overline Y$\footnote{Bản in ghi
$\overline f:\overline X\longrightarrow\overline X$; đối miền phải là
$\overline Y$, vì đây là cấu xạ được cảm sinh từ
$f:X\longrightarrow Y$.} của các lược đồ hình thức,
và cấu xạ này cũng proper. Nếu $F$ là một bó coherent trên $X$, thì
$\overline F$ là một bó coherent trên $\overline X$. Trong Định lý 1, ta quên
$X,Y,F$ và chỉ xét cấu xạ proper $\overline f$ của các lược đồ hình thức cùng
bó coherent $\overline F$ trên $\overline X$. (Tuy nhiên, báo cáo viên mới chỉ
viết một chứng minh đầy đủ cho trường hợp xuất phát từ $X,Y,f,F$.)

\result{ĐỊNH LÝ}{1 (định lý hữu hạn)}
\begin{enumerate}
  \item[\textit{i.}] Các $R^{q}\overline f_{*}(\overline F)$ là những bó
    coherent trên $\overline Y$.
  \item[\textit{ii.}] Các đồng cấu tự nhiên
    \[
      R^{q}\overline f_{*}(\overline F)
      \longrightarrow \varprojlim_n R^{q}f_{n*}(F_n)
    \]
    là những đẳng cấu.
\end{enumerate}

Trong mệnh đề này, ta giả sử đã chọn một bó con coherent $J$ của $\Osh_Y$ xác
định $Y'$; lấy nghịch ảnh, ta được một bó con coherent của $\Osh_X$ xác định
$X'$, và do đó xác định $F_n,X_n,Y_n$ cùng
$f_n:X_n\longrightarrow Y_n$ như trong mục số 2. Nếu xuất phát từ một cấu xạ
proper tùy ý giữa hai lược đồ hình thức, những thay đổi nhỏ cần thực hiện khi
viết tường minh các ký hiệu này là hiển nhiên.

Định lý 1 chỉ liên quan đến \fquote{đối đồng điều hình thức}. Định lý tiếp theo

\src{197}
liên hệ nó với \fquote{đối đồng điều đại số} và tương tự một định lý quen thuộc
của Serre [4] về sự so sánh giữa đối đồng điều đại số và đối đồng điều giải
tích.

\result{ĐỊNH LÝ}{2 (Định lý so sánh thứ nhất)} Các
$R^{q}f_{*}(F)$ là những bó coherent trên $Y$ (đây là một trường hợp riêng của
Định lý 1), và các đồng cấu tự nhiên
\[
  \overline{R^{q}f_{*}(F)}
  \longrightarrow \varprojlim_n R^{q}f_{n*}(F_n)
\]
là những đẳng cấu.

\result{HỆ QUẢ}{1} Ta có các đẳng cấu chính tắc:
\[
  \overline{R^{q}f_{*}(F)}=R^{q}\overline f_{*}(\overline F).
\]

Với $q=0$, hệ quả này là một khái quát của \fquote{định lý cơ bản về các hàm
chỉnh hình} của Zariski; từ đó ta sẽ suy ra một khái quát của
\fquote{định lý liên thông} của Zariski. Cũng cần lưu ý rằng, trong khi Định
lý 1 (\textit{ii}) là tầm thường với $q=0$, điều đó hoàn toàn không đúng đối
với Định lý 2 hay với phát biểu tương đương của nó (Hệ quả 1). Thật vậy, chứng
minh tiến hành bằng quy nạp lùi theo $q$ (nó là tầm thường khi $q$ lớn, vì khi
ấy cả hai vế đều bằng không), nên trường hợp $q=0$ xuất hiện ở bước cuối của
quy nạp và, có thể nói, là trường hợp \fquote{khó nhất}.

\result{HỆ QUẢ}{2} Giả sử $Y=\operatorname{Spec}(A)$ và $Y'$ được xác định bởi
iđêan $J$ của $A$. Khi đó, với mọi bó coherent $F$ trên $X$, các
$H^{q}(X,F)$ là những $A$-môđun hữu hạn sinh, và các hoàn thành $J$-adic của
chúng là các $H^{q}(\overline X,\overline F)$.

Cuối cùng, áp dụng hệ quả này cho $H=\Homsh_{\Osh_X}(F,G)$, ta thu được:

\result{HỆ QUẢ}{3} Giả sử $Y=\operatorname{Spec}(A)$ và $Y'$ được xác định bởi
iđêan $J$ của $A$. Cho $F,G$ là hai bó coherent trên $X$; khi đó
$\operatorname{Hom}(F,G)$ là một $A$-môđun hữu hạn sinh, và hoàn thành
$J$-adic của nó được đồng nhất với
$\operatorname{Hom}(\overline F,\overline G)$.

Dĩ nhiên, ánh xạ tự nhiên
$\operatorname{Hom}(F,G)\longrightarrow
\operatorname{Hom}(\overline F,\overline G)$ là ánh xạ gán cho một đồng cấu
$u:F\longrightarrow G$ phần mở rộng \fquote{theo liên tục}
$\overline u:\overline F\longrightarrow\overline G$ của nó (nhờ đó
$\overline F$ trở thành một hàm tử theo $F$).

Bây giờ giả sử $A$ tách và đầy đủ đối với tôpô $J$-adic; khi đó Hệ quả 2 và
Hệ quả 3 ở trên cho:
\[
  H^{q}(X,F)=H^{q}(\overline X,\overline F),
  \qquad
  \operatorname{Hom}(F,G)
  =\operatorname{Hom}(\overline F,\overline G).
\]

\src{198}
Đẳng thức cuối cùng này cho thấy phạm trù các bó coherent trên $X$ được đồng
nhất với một \emph{phạm trù con} (các cấu xạ là những cấu xạ cảm sinh) của
phạm trù các bó coherent trên $\overline X$. Thật ra, ta còn có kết quả mạnh
hơn sau đây:

\result{ĐỊNH LÝ}{3} Một bó môđun trên $\overline X$ là coherent khi và chỉ khi
nó đẳng cấu với một bó có dạng $\overline F$, trong đó $F$ là một bó coherent
trên $X$ (được xác định duy nhất đến đẳng cấu chính tắc theo Định lý 2, Hệ quả
3). [Nhắc lại rằng lúc này $Y=\operatorname{Spec}(A)$, với $A$ là một vành
tôpô Noether, tách và đầy đủ đối với tôpô $J$-adic.]

\result{HỆ QUẢ}{1} Các lược đồ con đóng của $X$ tương ứng song ánh với các
lược đồ con hình thức đóng của $\overline X$.

Thật vậy, chúng lần lượt tương ứng với các bó con coherent của $\Osh_X$ và
$\overline{\Osh}_X$. Xem đồ thị của các cấu xạ như những lược đồ con đóng, từ
Hệ quả 1 ta suy ra:

\result{HỆ QUẢ}{2} Cho $X,Z$ là hai lược đồ proper trên $A$ (một vành Noether
$J$-adic, tách và đầy đủ). Khi đó ánh xạ
$g\longmapsto\overline g$ xác định một tương ứng song ánh giữa các $Y$-cấu xạ
từ $X$ đến $Z$ và các $\overline Y$-cấu xạ từ $\overline X$ đến
$\overline Z$.

Nói cách khác, các lược đồ đại số proper trên $A$ tạo thành một
\emph{phạm trù con} (các cấu xạ là những cấu xạ cảm sinh) của phạm trù các
lược đồ hình thức proper trên $\overline Y$. Tuy nhiên, cần lưu ý rằng có những
lược đồ hình thức proper trên $\overline Y$ không \fquote{đại số hóa được},
nghĩa là không đẳng cấu với một $\overline X$ nào, trong đó $X$ proper trên
$A$ (cũng như có những đa tạp giải tích phức compact không phát sinh từ các đa
tạp đại số xác định trên trường số phức). Những lược đồ hình thức như vậy xuất
hiện một cách tự nhiên trong \fquote{lý thuyết moduli}. Dẫu vậy, ta ghi nhận
một trường hợp riêng đáng chú ý trong đó một lược đồ hình thức đại số hóa được:

\result{ĐỊNH LÝ}{4} Cho $A$ là một vành địa phương Noether đầy đủ có trường
thặng dư $k$, và cho $\mathfrak X$ là một lược đồ hình thức proper trên $A$
(được trang bị tôpô $\mathfrak m(A)$-adic). Giả sử:
\begin{enumerate}
  \item[\textit{i.}] các vành địa phương của $\Osh_{\mathfrak X}$ là những
    $A$-môđun phẳng; tương đương, nếu trang bị $\Osh_{\mathfrak X}$ và $A$
    các lọc xác định bởi các lũy thừa của iđêan cực đại của $A$, thì các đối
    tượng phân bậc liên kết thỏa mãn hệ thức
    \[
      \operatorname{gr}(\Osh_{\mathfrak X})
      \simeq
      \operatorname{gr}^{0}(\Osh_{\mathfrak X})
      \otimes_k\operatorname{gr}(A).
    \]
  \item[\textit{ii.}] khi xem
    $\mathfrak X_0=\mathfrak X\otimes_A k$ như một lược đồ đại số trên $k$,
    ta có
    \[
      H^2(\mathfrak X_0,\Osh_{\mathfrak X_0})=0.
    \]
  \item[\textit{iii.}] $\mathfrak X_0$ là xạ ảnh.
\end{enumerate}
Khi đó $\mathfrak X$ đại số hóa được; chính xác hơn, nó đẳng cấu với
$\overline X$, trong đó $X$ là một $A$-lược đồ xạ ảnh.

Các điều kiện \textit{(ii)} và \textit{(iii)} đặc biệt được thỏa mãn nếu
$\mathfrak X_0$ là một đường cong trơn trên $k$; vì thế Định lý 4 có thể được
áp dụng, chẳng hạn, trong \fquote{lý thuyết moduli} của các đường cong có
giống cho trước\ldots{} Ta phác họa cách chứng minh Định lý 4. Người ta chứng
minh (xem Mệnh đề 3 dưới đây) rằng \textit{(i)} và \textit{(ii)} kéo theo:
mọi bó coherent trên $\mathfrak X_0$ đẳng cấu địa phương với bó cấu trúc đều
nhận được bằng phép rút gọn từ một bó cùng loại trên $\mathfrak X$. Bắt đầu từ
một bó \fquote{ample} trên $\mathfrak X_0$ (tồn tại theo \textit{(iii)}), ta
nâng nó thành một bó khả nghịch trên $\mathfrak X$; rồi dùng Định lý 1 để chứng
minh rằng một lũy thừa thích hợp của bó ấy xác định một phép nhúng
$\mathfrak X$ vào hoàn thành hình thức của một lược đồ $\mathbf P_A^r$
(\fquote{không gian xạ ảnh chuẩn} chiều $r$ trên $A$).

Về chứng minh các Định lý 1 đến 3, xem [1].

\src{199}
\fsection{4}{Áp dụng vào định lý liên thông và \fquote{Định lý chính} của
Zariski}

Cho $f:X\longrightarrow Y$ là một cấu xạ proper của các lược đồ. Theo định lý
hữu hạn, $f_*(\Osh_X)=\mathcal A$ là một bó coherent trên $Y$; hơn nữa, đó là
một bó các đại số giao hoán, nên tương ứng với một $Y$-lược đồ
$g:Y'\longrightarrow Y$ hữu hạn trên $Y$ (được xác định bởi các điều kiện: nó
affine trên $Y$, nghĩa là nghịch ảnh của một tập mở affine là affine, và
$g_*(\Osh_{Y'})=\mathcal A$). Hiển nhiên khi đó $f$ phân tích chính tắc thành
$f=gf'$, trong đó $f':X\longrightarrow Y'$ là một cấu xạ từ $X$ đến $Y'$ thỏa
mãn $f'_*(\Osh_X)=\Osh_{Y'}$. Phân tích này của $f$ được gọi là
\emph{phân tích Stein} của $f$. Áp dụng Định lý so sánh thứ nhất, Hệ quả 1,
cho $f'$ và cho phần của $Y'$ rút gọn về một điểm $y'$, ta thấy
$f'^{-1}(y')=X'$ liên thông (nếu không, các tiết diện hình thức của $X$ dọc
theo $X'$ sẽ không tạo thành một vành địa phương, trong khi hoàn thành của
$f'_*(\Osh_X)_{y'}=\Osh_{Y',y'}$ lại là địa phương!). Ta đã chứng minh:

\result{ĐỊNH LÝ}{5 (\fquote{Định lý liên thông} của Zariski)} Cho
$f:X\longrightarrow Y$ là một cấu xạ proper. Khi đó $f$ phân tích duy nhất
(đến đẳng cấu) thành $f=gf'$, trong đó $g:Y'\longrightarrow Y$ là hữu hạn và
$f':X\longrightarrow Y'$ thỏa mãn
$f'_*(\Osh_X)=\Osh_{Y'}$ (do đó
$g_*(\Osh_{Y'})=f_*(\Osh_X)$). Các thớ của $f'$ liên thông; nói cách khác, tập
các thành phần liên thông của một thớ $f^{-1}(y)$ của $f$ tương ứng song ánh
với tập các điểm của $Y'$ nằm trên $y$, tức là với tập các iđêan cực đại của
$f_*(\Osh_X)_y$.

Các biến thể quen thuộc của định lý liên thông được suy ra ngay lập tức. Ở đây
ta chỉ ghi nhận hệ quả sau.

\src{200}
\result{HỆ QUẢ}{1} Một điểm $x$ của $X$ cô lập trong thớ $f^{-1}(y)$ khi và
chỉ khi thớ $f'^{-1}(y')$ (trong đó $y'=f'(x)$) rút gọn về $x$; tương đương,
$f'$ cảm sinh một đẳng cấu từ một lân cận của $x$ lên một lân cận của $y'$.
Tập các điểm như vậy là một tập mở $U$, và $f'$ cảm sinh một đẳng cấu từ $U$
lên một tập mở của $Y'$.

Để chứng minh $f'$ là một đẳng cấu địa phương tại $x$, ta lưu ý rằng $f'$ cảm
sinh một đẳng cấu $\Osh_{y'}\longrightarrow\Osh_x$, như suy ra từ
$f'_*(\Osh_X)=\Osh_{Y'}$ và từ sự kiện sau: vì $f'$ là một ánh xạ đóng có thớ
trên $y'$ rút gọn về $x$, các tập $f'^{-1}(V)$ chạy qua một hệ cơ bản các lân
cận của $x$ khi $V$ chạy qua một hệ cơ bản các lân cận của $y'$. Từ đó suy ra
ngay kết quả sau, do Chevalley chứng minh trong trường hợp
\fquote{hình học}.

\result{HỆ QUẢ}{2} Một cấu xạ $f$ là hữu hạn khi và chỉ khi nó proper và các
thớ của nó hữu hạn.

Thật vậy, trong trường hợp ấy $f'$ là một đẳng cấu theo kết quả vừa chứng minh.

Cho $f:X\longrightarrow Y$ là một cấu xạ không nhất thiết proper, nhưng giả sử
$X$ nằm như một tập mở trong một $Y$-lược đồ proper
$\widetilde f:\widetilde X\longrightarrow Y$ (điều này đặc biệt đúng nếu $f$
là quasi-xạ ảnh). Áp dụng Hệ quả 1, ta thấy $\widetilde f'$ cảm sinh một đẳng
cấu từ tập $U$ các điểm của $X$ cô lập trong thớ của chúng lên một tập mở của
$Y'$ (và $U$ quả thật là tập mở). Từ đó suy ra dạng toàn cục sau đây của
\fquote{Định lý chính} của Zariski.

\result{ĐỊNH LÝ}{6 (\fquote{Định lý chính} của Zariski)} Cho
$f:X\longrightarrow Y$ là một cấu xạ kiểu hữu hạn. Khi đó tập $U$ các điểm
của $X$ cô lập trong thớ của chúng là mở; nếu $f$ tách,\footnote{Bản đính chính
năm 1962 thay giả thiết in là \fquote{$f$ quasi-xạ ảnh} bằng giả thiết yếu hơn
\fquote{$f$ tách}, theo SGA VIII, 6.2: mọi cấu xạ quasi-hữu hạn và tách đều
quasi-xạ ảnh.} thì $U$ đẳng cấu trên $Y$ với một phần mở của một lược đồ $Y'$
hữu hạn trên $Y$.

Vì một cấu xạ kiểu hữu hạn là affine địa phương, và do đó \textit{a fortiori}
quasi-xạ ảnh địa phương, các dạng địa phương quen thuộc của Định lý chính được
suy ra ngay từ Định lý 6.

\src{201}
\result{MỆNH ĐỀ}{1} Cho $f:X\longrightarrow Y$ là một cấu xạ proper và $F$
là một bó trên $X$, coherent và phẳng trên $Y$. Cho $y\in Y$, cho $q$ là một
số nguyên, và giả sử $H^q(X_y,F_y)=0$. Khi đó $R^qf_*(F)$ triệt tiêu trong một
lân cận của $y$, và với mọi $n$, đồng cấu tự nhiên
\[
  R^{q-1}f_*(F)_y\longrightarrow
  H^{q-1}(X_y,F/\mathfrak m_y^nF)
\]
là toàn ánh.

Đặc biệt, nếu $f$ là một cấu xạ phẳng (nghĩa là $\Osh_X$ phẳng trên $Y$), thì
mọi bó coherent địa phương tự do $F$ trên $X$ đều phẳng trên $Y$. Cho $F,G$ là
hai bó như vậy; áp dụng Mệnh đề 1 cho $\Homsh_{\Osh_X}(F,G)$ với $q=1$, ta
được:

\result{ĐỊNH LÝ}{7} Cho $f$ là một cấu xạ proper và phẳng, $F,G$ là hai bó
coherent địa phương tự do trên $X$, $y\in Y$, và giả sử
\[
  H^1\bigl(X_y,\Homsh_{\Osh_{X_y}}(F_y,G_y)\bigr)=0,
\]
thì mọi đồng cấu $u_0:F_y\longrightarrow G_y$ đều được cảm sinh bởi một đồng
cấu
\[
  u:F|V\longrightarrow G|V,
  \qquad V=f^{-1}(U),
\]
trong đó $U$ là một lân cận của $y$.

\result{HỆ QUẢ}{} Nếu $u_0$ là một đẳng cấu (tương ứng, một đơn cấu, một toàn
cấu), thì $u$ cũng vậy khi $U$ đủ nhỏ.

Đặc biệt:

\src{202}
\result{HỆ QUẢ}{2} Cho $E_0$ là một bó coherent địa phương tự do trên $X_y$
thỏa mãn
\[
  H^1\bigl(X_y;\Homsh_{\Osh_{X_y}}(E_0,E_0)\bigr)=0,
\]
thì hai bó địa phương tự do mà các hạn chế xuống $X_y$ đều đẳng cấu
với $E_0$ sẽ đẳng cấu trong một lân cận của $X_y$.

Do đó:

\result{HỆ QUẢ}{3} Giả sử $H^1(X_y,\Osh_{X_y})=0$. Khi đó hai bó khả nghịch
trên $X$ (nghĩa là địa phương đẳng cấu với $\Osh_X$) mà các hạn chế xuống
$X_y$ đẳng cấu với nhau thì bản thân chúng cũng đẳng cấu trong một lân cận của $X_y$.\footnote{Bản in không nêu rõ lân cận. Ở đây khôi phục phạm vi địa phương của Hệ quả 2, cũng được sử dụng trong chứng minh Mệnh đề 2.}

Suy ra:

\result{MỆNH ĐỀ}{2} Cho $Y$ là một lược đồ liên thông và $E$ là một bó
coherent địa phương tự do có hạng ít nhất $2$ trên $Y$;\footnote{Bản in không
nêu cận dưới cho hạng. Giả thiết này là cần thiết: nếu $E$ có
hạng $1$, thì $\mathbf P(E)=Y$ và số nguyên $n$ trong phân tích dưới đây không
còn được xác định duy nhất. Hệ quả về các tự đẳng cấu phía sau kế thừa cùng
giả thiết.} xét phân thớ không gian xạ ảnh $X=\mathbf P(E)$ liên kết với $E$,
được trang bị bó khả nghịch quen thuộc $\Osh_X(1)$. Khi đó mọi bó khả nghịch
$L$ trên $X$ đều đẳng cấu với một bó có dạng
\[
  f^*(L')\otimes\Osh_X(n),
\]
trong đó $L'$ là một bó khả nghịch trên $Y$ và $n$ là một số nguyên. Số nguyên
vừa nêu được xác định duy nhất, còn $L'$ được xác định đến đẳng cấu.

Hệ quả 3 ở trên chứng minh rằng $L$ đẳng cấu với một $\Osh_X(n)$ nào đó trong
một lân cận của mỗi thớ. Phần còn lại gần như chỉ là hình thức.

Mệnh đề 2 cho phép xác định các $Y$-cấu xạ từ $X=\mathbf P(E)$ đến một phân
thớ xạ ảnh khác. Đặc biệt, ta được:

\result{HỆ QUẢ}{} Cho $u$ là một tự đẳng cấu của $X=\mathbf P(E)$. Khi đó tồn
tại một bó khả nghịch $L'$ trên $Y$ và một đẳng cấu $v$ từ $E$ lên
$E\otimes L'$ sao cho $u$ là đẳng cấu tương ứng
\[
  \mathbf P(E)\xrightarrow{\sim}\mathbf P(E\otimes L')=\mathbf P(E);
\]
cặp $(v,L')$ được xác định đến đẳng cấu.

Cho $\Gamma$ là tập các lớp của những bó khả nghịch $L'$ trên $Y$ sao cho
$E\otimes L'$ đẳng cấu với $E$. Các phần tử của tập này đều là phần tử xoắn:
thật vậy, nếu $n$ là hạng của $E$, thì lấy lũy thừa ngoài bậc $n$ cho thấy
nhất thiết phải có
$L'^{\otimes n}\simeq\Osh_Y$. Vì thế, hệ quả có thể được phát biểu bằng dãy
khớp các nhóm
\[
  e\longrightarrow
  \operatorname{Aut}(E)/\Gamma(Y,\Osh_Y^*)
  \longrightarrow \operatorname{Aut}_Y(X)
  \longrightarrow \Gamma\longrightarrow e
\]
(dãy này cũng suy ra từ dãy khớp đối đồng điều phát sinh từ dãy khớp sau đây
của các bó nhóm
\[
  e\longrightarrow \Osh_Y^*
  \longrightarrow \underline{\operatorname{Aut}}(E)
  \longrightarrow \underline{\operatorname{Aut}}_Y(X)
  \longrightarrow e,
\]
trong đó $\Osh_Y^*$ là bó các \fquote{đơn vị} của $\Osh_Y$, được đồng nhất với
tâm của $\underline{\operatorname{Aut}}(E)$).

\src{203}
\fsection{6}{Ứng dụng vào các định lý tồn tại và duy nhất cho bó và lược đồ
trên một vành đầy đủ theo tôpô \texorpdfstring{$J$}{J}-adic}

Định lý 7 đã cho một kết quả duy nhất đối với các bó coherent địa phương tự do,
bằng cách dùng các Định lý 1 và 2. Giờ đây, nhờ Định lý 3, ta suy ra các định
lý tồn tại cho bó, cấu xạ lược đồ hoặc lược đồ. Trong phần còn lại, $A$ ký hiệu
một vành Noether địa phương, tách và đầy đủ. Phương pháp tổng quát luôn là thực
hiện các phép dựng hình thức, về bản chất tức là làm hình học đại số trên một
vành Artin, rồi dùng ba định lý cơ bản để rút ra các kết luận có tính chất
“đại số”.

\result{MỆNH ĐỀ}{3} Cho $\mathfrak X$ là một lược đồ hình thức proper và phẳng
trên $A$, và cho $F_0$ là một bó địa phương tự do trên $X_0$ sao cho
\[
  H^2\bigl(X_0,\Homsh_{\Osh_{X_0}}(F_0,F_0)\bigr)=0.
\]
Khi đó tồn tại một bó địa phương tự do $F$ trên $\mathfrak X$ mà hạn chế xuống
$X_0$ đẳng cấu với $F_0$. (Hơn nữa, $F$ này duy nhất đến đẳng cấu nếu
$H^1(X_0,\Homsh_{\Osh_{X_0}}(F_0,F_0))=0$.)

Ta lần lượt dựng các bó địa phương tự do $F_n$ trên các $X_n$, mỗi bó cảm sinh
bó đứng trước. Cản trở đối với việc dựng $F_n$ nằm trong
\[
  H^2\bigl(X_0,\Homsh_{\Osh_{X_0}}(F_0,F_0)\bigr)
  \otimes_{A/J}(J^n/J^{n+1}),
\]
nên bằng không theo giả thiết. Bây giờ áp dụng Định lý 3, ta được:

\result{HỆ QUẢ}{1} Cho $X$ là một lược đồ proper và phẳng trên $A$, và cho
$F_0$ như trên. Khi đó tồn tại một bó địa phương tự do $F$ trên $X$ mà hạn chế
xuống $X_0$ đẳng cấu với $F_0$. Hơn nữa, $F$ này duy nhất đến đẳng cấu nếu
$H^1(X_0,\Homsh_{\Osh_{X_0}}(F_0,F_0))=0$.

Cho $X_0$ là một lược đồ kiểu hữu hạn trên trường $k$; giả sử $X_0$ trơn (ở
đây hiểu là trơn tuyệt đối) trên $k$, nhưng không nhất thiết proper trên $k$.
Cho $A$ là một vành Artin địa phương có trường thặng dư $k$. Ta muốn tìm các
lược đồ $X$ phẳng trên $A$ sao cho $X\otimes_A k=X_0$ (đây là điểm xuất phát
của “lý thuyết moduli” hoặc của “các biến thiên cấu trúc” của $X_0$). Việc cho
một $X$ như vậy tương đương với việc cho, trên không gian tôpô $X_0$, một bó
$\Osh_X$ được trang bị các cấu trúc sau:
\begin{enumerate}
  \item[\textit{i.}] Nó là một bó các $A$-đại số.
  \item[\textit{ii.}] Nó được trang bị một đồng cấu tăng cường
    $\Osh_X\longrightarrow\Osh_{X_0}$ (tương thích với các cấu trúc
    $A$-đại số); các dữ liệu này phải thỏa các điều kiện sau: phép tăng cường
    cảm sinh một đẳng cấu
    $\Osh_X\otimes_A k\xrightarrow{\sim}\Osh_{X_0}$; $\Osh_X$ phẳng trên
    $A$, nghĩa là đối tượng phân bậc liên kết với $\Osh_X$, được lọc bởi các
    lũy thừa của iđêan cực đại $\mathfrak m$ của $A$, đẳng cấu với
    $\operatorname{gr}^0(\Osh_X)\otimes_k\operatorname{gr}(A)$; tương đương,
    ta có các đẳng cấu
    \[
      \mathfrak m^n\Osh_X/\mathfrak m^{n+1}\Osh_X
      =\Osh_{X_0}\otimes_k(\mathfrak m^n/\mathfrak m^{n+1}).
    \]
\end{enumerate}
Sự kiện cơ bản là như sau:

\src{204}
\result{ĐỊNH LÝ}{8} Cho $X_0$ là một lược đồ kiểu hữu hạn và trơn trên trường
$k$, và giả sử $X_0$ affine. Cho $A$ là một vành Artin địa phương có trường
thặng dư $k$. Khi đó tồn tại một $A$-lược đồ $X$ phẳng trên $A$ sao cho
$X\otimes_A k=X_0$, và hai lược đồ bất kỳ như vậy tất yếu đẳng cấu.

Lưu ý rằng đẳng cấu nói trên không chính tắc, vì nói chung $X$ sẽ có những tự
đẳng cấu trên $A$ không tầm thường cảm sinh đồng nhất trên $X_0$. Mặt khác,
nói chung không có cách chọn “chính tắc” một $X$ thỏa các điều kiện đã nêu,
trừ khi $A$ là một $k$-đại số (trường hợp cùng đặc số); khi ấy có thể lấy
$X=X_0\otimes_k A$, nghĩa là $\Osh_X=\Osh_{X_0}\otimes_k A$ (dù $X_0$ affine
hay không). Trong trường hợp khác đặc số, khi $X_0$ không affine, nói chung ta
không biết liệu có thể “nâng” $X_0$ thành một $X$ xác định trên $A$ hay không.
Tuy nhiên, cho $n$ là một số nguyên $>0$, đặt
$A_{n-1}=A/\mathfrak m^n$, và giả sử đã nâng $X_0$ thành một
$A_{n-1}$-lược đồ phẳng $X_{n-1}$; ta muốn nâng $X_{n-1}$ thành một
$A_n$-lược đồ phẳng $X_n$. Theo Định lý 8, ta đã biết điều này thực hiện được
ở địa phương. Hơn nữa, dễ kiểm tra rằng nếu $U_n$ nâng một tập mở $U_{n-1}$
của $X_{n-1}$, thì bó nhóm các tự đẳng cấu của $U_n$ (cảm sinh đồng nhất trên
$U_{n-1}$) đẳng cấu chính tắc với
\[
  \Theta_{X_0/k}\otimes_k(\mathfrak m^n/\mathfrak m^{n+1})
\]
được hạn chế xuống $U_{n-1}$, và đặc biệt có tính giao hoán
($\Theta_{X_0/k}$ ký hiệu bó các mầm của các $k$-đạo hàm trên $X_0$). Từ đó
dễ suy ra rằng có một cản trở đối với việc dựng $X_n$ nâng $X_{n-1}$, nằm
trong
\[
  H^2(X_0,\Theta_{X_0/k})\otimes_k
  \mathfrak m^n/\mathfrak m^{n+1}.
\]
Do đó:

\result{HỆ QUẢ}{1} Cho $X_0$ là một lược đồ kiểu hữu hạn trơn trên $k$, và
giả sử $H^2(X_0,\Theta_{X_0/k})=0$. Khi đó, với mọi vành Artin địa phương
$A$ có trường thặng dư $k$, tồn tại một $A$-lược đồ phẳng $X$ sao
cho $X\otimes_A k=X_0$.

Hơn nữa, nếu đã tìm được một $X$ phẳng trên $A$ nâng $X_0$, thì theo Định lý
8, tập các lớp (đến đẳng cấu) của các $A$-lược đồ phẳng nâng $X_0$ được đồng
nhất với
\[
  H^1\bigl(X_0,\underline{\operatorname{Aut}}(X)\bigr),
\]
trong đó, dĩ nhiên, $\underline{\operatorname{Aut}}(X)$ ký hiệu bó các mầm
của những tự đẳng cấu của bó $A$-đại số $\Osh_X$ tương thích với phép tăng
cường. Lọc của $\Osh_X$ xác định một lọc của
$\underline{\operatorname{Aut}}(X)$; thương của bó này cho nhóm con thứ
$n$ trong lọc là $\underline{\operatorname{Aut}}(X_n)$. Đối tượng phân bậc
liên kết với lọc này có tính giao hoán và được đồng nhất với
$\Theta_{X_0/k}\otimes_k\operatorname{gr}(A)$. Đặc biệt, nếu
$\mathfrak m^{n+1}$ là lũy thừa đầu tiên của $\mathfrak m$ bằng không, thì
$F^n(\underline{\operatorname{Aut}}(X))$ (bậc cuối của lọc) nằm trong tâm của
$\underline{\operatorname{Aut}}(X)$ và đẳng cấu với
$\Theta_{X_0/k}\otimes_k\mathfrak m^n$; nó cũng là bó các mầm của những tự
đẳng cấu của $X$ cảm sinh đồng nhất trên
$X_{n-1}=X\otimes_A A/\mathfrak m^n$. Dùng các kết quả này, ta lập tức thu
được các mệnh đề sau:

\src{205}
\result{HỆ QUẢ}{2} Cho $X_0$ là một lược đồ kiểu hữu hạn trơn trên $k$, cho
$A$ là một vành Artin địa phương có trường thặng dư $k$ và iđêan cực đại
$\mathfrak m$, giả sử $\mathfrak m^{n+1}=0$, đặt
$A_{n-1}=A/\mathfrak m^n$, và cuối cùng cho $X_{n-1}$ là một
$A_{n-1}$-lược đồ phẳng sao cho $X_{n-1}\otimes_A k=X_0$. Khi đó tập các
lớp (đến một đẳng cấu cảm sinh đồng nhất trên $X_{n-1}$) của các $A$-lược đồ
phẳng $X_n$ sao cho $X_n\otimes_A A_{n-1}=X_{n-1}$ hoặc rỗng, hoặc là một
không gian thuần nhất chính dưới
\[
  H^1(X_0,\Theta_{X_0/k})\otimes_k\mathfrak m^n.
\]

(Lưu ý rằng nói chung không có điểm gốc ưu tiên trong không gian vừa nêu, vì
không có cách ưu tiên nào để nâng $X_{n-1}$ thành $X_n$.)

\result{HỆ QUẢ}{3} Cho $X_0$ là một lược đồ kiểu hữu hạn trơn trên $k$, và
giả sử $H^1(X_0,\Theta_{X_0/k})=0$. Khi đó, với mọi vành Artin địa phương
$A$ có trường thặng dư $k$, tồn tại nhiều nhất một (đến đẳng cấu) $A$-lược đồ
phẳng $X$ sao cho $X\otimes_A k=X_0$.

Các Hệ quả 1 và 3 lập tức kéo theo những mệnh đề bề ngoài tổng quát hơn, thu
được khi chỉ giả sử $A$ là một vành Noether địa phương đầy đủ có trường thặng
dư $k$, miễn là ta đưa $X$ vào dưới dạng một lược đồ hình thức trên $A$:

\result{ĐỊNH LÝ}{9} Cho $k$ là một trường và $X_0$ là một lược đồ kiểu hữu
hạn trơn trên $k$. Với mỗi vành Noether địa phương đầy đủ $A$ có trường thặng
dư $k$, gọi $F(A)$ là tập các lớp (đến một đẳng cấu cảm sinh đồng nhất trên
$X_0$) của các lược đồ hình thức $\mathfrak X$ trên $A$, kiểu hữu hạn và phẳng
trên $A$, sao cho $\mathfrak X\otimes_A k=X_0$. Với các ký hiệu ấy:
\begin{enumerate}
  \item[\textit{i.}] Nếu $H^1(X_0,\Theta_{X_0/k})=0$, thì với mọi $A$,
    $F(A)$ có nhiều nhất một phần tử.
  \item[\textit{ii.}] Nếu $H^2(X_0,\Theta_{X_0/k})=0$, thì với mọi $A$,
    $F(A)$ có ít nhất một phần tử.
\end{enumerate}

\src{206}
\result{HỆ QUẢ}{1} Giả sử $X_0$ proper trên $k$. Dưới điều kiện
\textit{(i)}, với mọi $A$ tồn tại nhiều nhất một (đến một đẳng cấu cảm sinh
đồng nhất trên $X_0$) lược đồ $X$ proper và phẳng trên $A$ sao cho
$X\otimes_A k=X_0$.

Ở đây ta dùng Định lý 3, Hệ quả 2. Chẳng hạn:

\result{HỆ QUẢ}{2} Nếu $X$ là một lược đồ proper và phẳng trên $A$ sao cho
$X\otimes_A k$ đẳng cấu với lược đồ xạ ảnh chuẩn $\mathbf P_k^r$ chiều $r$
trên $k$, thì $X$ đẳng cấu với $\mathbf P_A^r$.

(Kết quả này cũng có thể suy ra từ Mệnh đề 3, Hệ quả 1.)

\result{HỆ QUẢ}{3} Cho $X_0$ là một lược đồ xạ ảnh trơn trên $k$, và giả sử
\[
  H^2(X_0,\Osh_{X_0})=H^2(X_0,\Theta_{X_0/k})=0.
\]
Khi đó, với mọi $A$, tồn tại một lược đồ xạ ảnh $X$, phẳng trên $A$, sao cho
$X\otimes_A k=X_0$.

Kết hợp Định lý 9\textit{(ii)} với Định lý 4. Đặc biệt:

\result{HỆ QUẢ}{4} Cho $X_0$ là lược đồ của một đường cong đại số đầy đủ và
trơn trên $k$. Khi đó, với mọi vành Noether địa phương đầy đủ $A$ có trường
thặng dư $k$, tồn tại một “lược đồ đường cong trơn” $X$ trên $A$ sao cho
$X\otimes_A k=X_0$.

\result{NHẬN XÉT}{}

1) Các Hệ quả 3 và 4 đặc biệt đáng chú ý khi $k$ có đặc số $p\ne0$: ta lấy
$A$ là một vành định giá rời rạc có đặc số $0$ và trường thặng dư $k$; chẳng
hạn, “$A$ nhỏ nhất có thể”, tức là vành có iđêan cực đại sinh bởi $p$. (Thật
vậy, theo các định lý của Cohen, chỉ cần biết các Hệ quả 3 và 4 cho một vành
$A$ như vậy.) Về điểm này, cần lưu ý rằng theo các chuyên gia, chưa biết có
tồn tại những lược đồ trên một trường $k$ không phải là phép rút gọn modulo
$p$ của một lược đồ phẳng xác định trên một vành $A$ như vậy hay không. Ít
nhất, các kết quả của mục này cung cấp một phương pháp để khảo sát câu hỏi ấy
một cách có hệ thống. Trước hết nên xét xem cản trở đầu tiên xuất hiện trong
$H^2(X_0,\Theta_{X_0/k})$ có nhất thiết bằng không hay không.

\emph{Bổ sung năm 1962.} --- J.-P. Serre đã dựng một đa tạp xạ ảnh trơn chiều
3 trên một trường đóng đại số $k$ có đặc số $p>0$, không thu được bằng phép
rút gọn từ một lược đồ proper trên một vành địa phương nguyên có trường thặng
dư $k$ và trường phân thức đặc số không. Nghe nói Mumford đã tìm được một ví
dụ tương tự với một mặt xạ ảnh trơn.

2) Cần lưu ý rằng Định lý 3 và kỹ thuật tương ứng chỉ áp dụng khi vành cơ sở
(địa phương, để cố định ý tưởng) là đầy đủ.

Muốn chuyển từ các kết quả đã biết cho vành đầy đủ hóa của một vành địa phương
sang các kết quả tương ứng cho chính vành địa phương ấy, ta cần một “định lý
cơ bản” thứ tư, mà phát biểu dứt khoát của nó vẫn còn phải tìm.

\src{207}
3) Hãy so sánh các kết quả của mục này (đặc biệt là các Hệ quả 1 và 2 ở trên)
và của mục kế tiếp với các kết quả của Kodaira--Spencer về sự biến thiên của
các cấu trúc phức. Dùng định lý phỏng đoán vừa nói đến, dưới các điều kiện của
Hệ quả 1 nhưng không còn giả sử $A$ đầy đủ, ta hẳn có thể kết luận rằng tồn
tại một vành $A'$ chứa $A$, hữu hạn và không phân nhánh trên $A$, sao cho
$X\otimes_A A'$ và $X'\otimes_A A'$ đẳng cấu trên $A'$ (ở đây $X,X'$ là hai
$A$-lược đồ proper và phẳng đã cho, thỏa
$X\otimes_A k=X'\otimes_A k=X_0$). Ít nhất, điều này có thể chứng minh khi
$X_0=\mathbf P_k^r$, bằng cách dùng hệ quả của Mệnh đề 2. Trong mọi trường
hợp, khi $H^1(X_0,\Theta_{X_0/k})=0$, có thể chứng minh rằng các thớ của
$X,X'$ trên mọi điểm $y$ của $Y=\operatorname{Spec}(A)$ là đẳng cấu, ít nhất
sau khi chuyển sang bao đóng đại số của trường thặng dư $\kappa(y)$. (Có một
kết quả địa phương, dường như mạnh hơn, trong đó không còn cần giả sử $A$ là
địa phương.) Đối với “các biến thiên cấu trúc” của không gian xạ ảnh, ta cũng
nêu câu hỏi sau đây, do một bài toán tương ứng của Kodaira--Spencer gợi ra.
Cho $X$ là một lược đồ proper và phẳng trên vành địa phương nguyên $A$, có
trường phân thức $K$ và trường thặng dư $k$, và giả sử $X\otimes_A K$ đẳng
cấu với $\mathbf P_K^r$; liệu $X\otimes_A k=X_0$ có đẳng cấu với
$\mathbf P_k^r$ hay không (ít nhất trên bao đóng đại số của $k$)? Trong bài
toán này, có thể giả sử $A$ là một vành định giá rời rạc đầy đủ. Một bài toán
tương tự được đặt ra khi $X_0$ là một đa tạp abel.

\emph{Bổ sung năm 1962.} --- Khi ấy diễn giả đã bớt lạc quan hơn về các kết
quả phỏng đoán ở trên. Tuy nhiên, Hironaka đã giải bài toán về các biến thiên
cấu trúc của không gian xạ ảnh theo hướng khẳng định, còn Koizumi giải bài
toán tương tự cho các đa tạp abel.

\fsection{7}{Ứng dụng vào “lý thuyết moduli”}

Vì bản thân diễn giả mới chỉ vừa bắt đầu nghiên cứu lý thuyết này, ta buộc
phải giới hạn ở một vài chỉ dẫn. Để đơn giản, ta làm việc trên một trường $k$,
nghĩa là trong trường hợp cùng đặc số, mặc dù Định lý 8 dường như cũng cho phép
xử lý trường hợp tổng quát mà không có thay đổi bản chất nào. Hiện tại ta chưa
vượt quá giai đoạn “hình thức”; tuy nhiên, diễn giả hy vọng từ đó có thể dựng
được các lược đồ moduli thực sự trong một số trường hợp và, đặc biệt, với mỗi
số nguyên $g$, dựng một lược đồ trên các số nguyên đóng vai trò lược đồ moduli
phổ dụng cho các đường cong trơn có giống $g$.

\emph{Bổ sung năm 1962.} --- Mumford đã dựng các lược đồ moduli cho các đường
cong giống $g$ (xem SHMI). Hơn nữa, Định lý 10 chứng minh rằng các “lược đồ
Jacobi mức $n$” của lý thuyết moduli là không kỳ dị, và thậm chí trơn trên
$\mathbf Z$.

\src{208}
Trở lại tình huống và ký hiệu của Định lý 9, giờ đây giả sử $A$ là một đại số
địa phương hữu hạn chiều trên $k$, và để đơn giản trường này được giả sử đóng đại số.
Khi đó có thể xem $F(A)$ như một hàm tử hiệp biến theo $A$, nhận giá trị trong
phạm trù các tập hợp: một đồng cấu các $k$-đại số $A\longrightarrow B$ xác
định một ánh xạ $F(A)\longrightarrow F(B)$, vì mỗi $A$-lược đồ phẳng $X$ sao
cho $X\otimes_A k=X_0$ cho một $B$-lược đồ $X\otimes_A B$ có cùng các tính
chất. Giả sử có thể tìm được một $k$-đại số Noether địa phương đầy đủ
$\mathcal O$ và một đẳng cấu có tính hàm tử
\begin{equation}
  \tag{*}
  \operatorname{Hom}(\mathcal O,A)\xrightarrow{\sim}F(A)
\end{equation}
(trong đó vế thứ nhất ký hiệu các đồng cấu $k$-đại số). Dễ thấy một
$\mathcal O$ như vậy được xác định đến đẳng cấu chính tắc; khi ấy, ta gọi phổ
hình thức $\mathfrak Y$ của $\mathcal O$ (nghĩa là không gian tôpô chỉ gồm
một điểm và được trang bị một bó các vành tôpô thu về $\mathcal O$) là
\emph{lược đồ moduli hình thức} của $X_0$. (Lược đồ này không nhất thiết tồn
tại.) Cho $\mathfrak r$ là iđêan cực đại của $\mathcal O$; với mỗi $n$, đặt
$\mathcal O_n=\mathcal O/\mathfrak r^{n+1}$ (do đó $\mathcal O_0=k$). Khi
ấy, đồng cấu chính tắc $\mathcal O\longrightarrow\mathcal O_n$ là một phần tử
của $\operatorname{Hom}(\mathcal O,\mathcal O_n)$, nên xác định một phần tử
của $F(\mathcal O_n)$, tức là một $\mathcal O_n$-lược đồ phẳng $X_n$ có phép
rút gọn modulo $\mathfrak r$ bằng $X_0$. Các $X_n$ này suy ra từ nhau bằng
mở rộng vô hướng (ở đây tức là bằng các phép rút gọn), do đó chúng đến từ một
lược đồ hình thức $\mathfrak X$ xác định duy nhất trên lược đồ moduli hình
thức $\mathfrak Y$; $\mathfrak X$ phẳng trên $\mathfrak Y$ và
$\mathfrak X_0=X_0$. Khi đó, như thấy ngay, đẳng cấu (*) có được bằng cách
gắn với mỗi đồng cấu $k$-đại số $\mathcal O\longrightarrow A$ lớp của
$A$-lược đồ $\mathfrak X\otimes_{\mathcal O}A$ (nghĩa là, với mỗi cấu xạ các
$k$-lược đồ $Y'=\operatorname{Spec}(A)\longrightarrow\mathfrak Y$, ta gắn
$Y'$-lược đồ $\mathfrak X\otimes_{\mathfrak Y}Y'$ thu được bằng đổi cơ sở).
Hơn nữa, đẳng cấu (*) và mô tả vừa nêu vẫn còn đúng nếu chỉ giả sử $A$ là một
$k$-đại số Noether địa phương đầy đủ (không nhất thiết Artin). Dĩ nhiên, như
thường lệ, $\mathcal O$ hoàn toàn có thể có các phần tử lũy linh
\textit{a priori}, và dường như có thể tồn tại những trường hợp trong đó
$\mathcal O$ tự nó là Artin nhưng không đồng nhất với $k$. Điều này cho thấy
quan điểm Kodaira--Spencer (chỉ xét các $A$ là vành chính quy) không thích hợp
\textit{a priori} đến mức nào trong trường hợp tổng quát.

\src{209}
Ta còn phải nêu các điều kiện đủ để tồn tại một lược đồ moduli hình thức cho
$X_0$, được giả sử proper trên $k$. Nói chung, dễ nêu các điều kiện cần và đủ
đơn giản trên một hàm tử $A\longmapsto F(A)$ (từ các $k$-đại số địa phương
hữu hạn chiều đến các tập hợp) để nó có thể viết dưới dạng
$\operatorname{Hom}(\mathcal O,A)$ với một $\mathcal O$ thích hợp. Ta sẽ
không trình bày chi tiết các điều kiện ấy ở đây. Chỉ xin lưu ý rằng, trong
trường hợp đang xét, chúng áp đặt lên $X_0$ những điều kiện đối đồng điều không
tầm thường, và dường như khó có khả năng chúng luôn được thỏa mãn, dù diễn giả
chưa dựng được phản ví dụ nào. Trái lại, có vẻ hợp lý khi cho rằng điều kiện
$H^0(X_0,\Theta_{X_0/k})=0$ là điều kiện đủ (dù tuyệt nhiên không cần) cho sự
tồn tại của một lược đồ moduli hình thức. Ở đây ta chỉ phát biểu một định lý
trong một trường hợp đặc biệt đơn giản (định lý tương tự trong lý thuyết không
gian giải tích đã được biết rõ; xem Kodaira--Spencer), và định lý này được suy
ra không khó khăn từ các kết quả của mục trước:

\result{ĐỊNH LÝ}{10} Cho $X_0$ là một lược đồ proper và trơn trên trường $k$
sao cho
\[
  H^0(X_0,\Theta_{X_0/k})=H^2(X_0,\Theta_{X_0/k})=0.
\]
Khi đó tồn tại một lược đồ moduli hình thức cho $X_0$, tương ứng với một vành
địa phương chính quy $\mathcal O$ (nghĩa là một đại số chuỗi lũy thừa hình thức
trên $k$).

Như đã lưu ý, nói chung không đúng rằng lược đồ hình thức $\mathfrak X$ trên
$\mathcal O$ là đại số hóa được; tuy nhiên, ta biết điều đó đúng khi $X_0$ là
xạ ảnh và $H^2(X_0,\Osh_{X_0})=0$ (Định lý 4), chẳng hạn khi $X_0$ có chiều
1. Điều này đem lại một phần cơ sở cho hy vọng dựng được một lược đồ moduli
trên các số nguyên cho các đường cong có giống cho trước\ldots{}

Cũng cần lưu ý rằng các phương pháp như những phương pháp trình bày trong mục
này có thể áp dụng vào việc dựng và nghiên cứu các đa tạp Picard, cũng như vào
nhiều phép dựng khác. Ta sẽ sớm trở lại vấn đề này.

\fsection{8}{Ứng dụng vào nhóm cơ bản}

Các kỹ thuật đã trình bày cho phép bắt đầu nghiên cứu có hệ thống nhóm cơ bản
theo mô hình của lý thuyết tôpô. Hai định lý đầu tiên được phát biểu trong mục
này là những khái quát hóa các kết quả trong một công trình gần đây của
Lang--Serre.

Cho $X$ là một lược đồ. Một $X$-lược đồ $X'$ được gọi là một phủ không phân
nhánh của $X$ nếu
\begin{enumerate}
  \item[\textit{i.}] $X'$ hữu hạn trên $X$, nghĩa là nó được xác định bởi một
    bó coherent các đại số $\mathcal A=\mathcal A(X')$ trên $X$;
  \item[\textit{ii.}] $\mathcal A$ là một bó địa phương tự do trên $X$;
  \item[\textit{iii.}] Với mọi $x\in X$, thương
    \[
      \mathcal A_x/\mathfrak m_x\mathcal A_x
      =\mathcal A_x\otimes_{\Osh_x}\kappa(x)
    \]
    là một đại số khả phân trên $\kappa(x)$.
\end{enumerate}

\src{210}
Khái niệm phủ không phân nhánh này (do Serre và diễn giả đề xuất) có mọi tính
chất sơ cấp mà ta có thể kỳ vọng một cách hợp lý; ta sẽ không liệt kê chúng.
Chỉ xin nói rằng nó dẫn đến một lý thuyết Galois theo mô hình của lý thuyết
Galois cổ điển (và bao hàm lý thuyết ấy; các chứng minh thậm chí còn đơn giản
hơn khá nhiều so với những chứng minh thường được chấp nhận cho lý thuyết sau) và của lý thuyết
Galois về các phủ tôpô. Cụ thể hơn, ta gọi một \emph{điểm hình học} của một
lược đồ $X$ là một cấu xạ $a$ từ phổ $\xi$ của một trường đóng đại số $\Omega$
vào $X$, nghĩa là dữ liệu của một mở rộng đóng đại số của trường thặng dư
$\kappa(x)$ của một điểm $x=|a|$ thuộc $X$ (được gọi là \emph{vị trí} của
điểm hình học $a$). Nếu $X'$ là một phủ không phân nhánh của $X$, ta có thể
gắn với nó tập $E_a(X')$ các “điểm hình học của $X'$ nằm trên $a$”, nghĩa là
tập các cặp gồm một $x'\in X'$ nằm trên $x$ và một $\kappa(x)$-đồng cấu
$\kappa(x')\longrightarrow\Omega$. Như vậy (với $(X,a)$ cố định), ta thu được
một hàm tử $F(X,a)$ từ phạm trù $R(X)$ các phủ không phân nhánh $X'$ của $X$
đến phạm trù các tập hữu hạn. Nếu $X$ liên thông, cặp gồm $R(X)$ và $F(X,a)$
có các tính chất hình thức cần thiết để đẳng cấu với cặp tương tự được xác định
bởi một nhóm tôpô compact hoàn toàn không liên thông $\pi$ thích hợp (nghĩa là
một giới hạn xạ ảnh của các nhóm hữu hạn): ta lấy phạm trù $C(\pi)$ của các tập
hữu hạn $E$ trên đó $\pi$ tác động liên tục, cùng hàm tử quên
$F(\pi)(E)=E$ từ phạm trù này đến phạm trù các tập hữu hạn. Hơn nữa, nhóm
$\pi$ được xác định đến đẳng cấu chính tắc bởi điều kiện $(C(\pi),F(\pi))$
đẳng cấu với một cặp đã cho. Trong trường hợp này, $\pi$ được gọi là
\emph{nhóm cơ bản} của lược đồ liên thông $X$ tại điểm hình học $a$, và được
ký hiệu $\pi_1(X,a)$. Nếu $X$ không liên thông, ta thay nó bằng thành phần liên
thông chứa $x=|a|$. Ngược lại, nếu $X$ liên thông, thì các nhóm $\pi_1(X,a)$
ứng với hai điểm hình học $a',a''$ của $X$ là đẳng cấu (đẳng cấu ấy chỉ được
xác định đến các tự đẳng cấu trong), nên như thường lệ, ta có thể chọn $a$
theo cách thuận tiện nhất, chẳng hạn tại điểm tổng quát của $X$, với giả thiết
lược đồ này bất khả quy. Dĩ nhiên, $\pi_1(X,a)$ là một hàm tử hiệp biến theo lược
đồ có điểm gốc $(X,a)$. Khi đó, mọi phát biểu về việc phân loại các phủ không
phân nhánh\footnote{Bản in ghi “các phủ bất khả phân”; đoạn văn đang xét phạm
trù $R(X)$ của các phủ không phân nhánh và sự tương đương của nó với các tập
hữu hạn được trang bị một tác động liên tục của $\pi_1(X,a)$.} đều có thể
chuyển sang ngôn ngữ lý thuyết nhóm theo từ điển quen thuộc (ngoại trừ việc
phải lưu ý rằng ở đây ta làm việc với các nhóm tôpô).

\src{211}
Mục đích của ta là phát triển một tương tự của dãy đồng luân chính xác của các
không gian phân thớ đối với một cấu xạ proper $f:X\longrightarrow Y$. Hiển
nhiên, vì chưa biết các nhóm đồng luân bậc cao phải là gì, ta chỉ có thể thu
được những kết quả tất yếu chưa đầy đủ. Để áp dụng các định lý cơ bản của
mục n\textsuperscript{o} 3, trước hết ta phải phát biểu tường minh vài bổ đề sơ cấp về các lược đồ
trên trường hoặc trên vành Artin (đúng theo phương pháp tổng quát!).

\result{BỔ ĐỀ}{1} Cho $(X',a')$ là một phủ không phân nhánh có điểm gốc, liên
kết với một biểu diễn có điểm gốc của $\pi_1(X,a)$ trên một tập hữu hạn $E$
(có điểm được đánh dấu $e$). Khi đó cấu xạ chính tắc
$\pi_1(X',a')\longrightarrow\pi_1(X,a)$ đồng nhất nhóm thứ nhất với nhóm ổn
định của $e$ trong $\pi_1(X,a)$ (và do đó là đơn ánh).

\result{BỔ ĐỀ}{2} Cho $X$ là một lược đồ đại số trên trường $k$, và cho $k'$
là một mở rộng thuần bất khả phân của $k$. Khi đó mọi phủ không phân nhánh của
$X\otimes_k k'$ đều thu được bằng ảnh ngược (tức là bằng mở rộng vô hướng) từ
một phủ không phân nhánh của $X$, được xác định duy nhất đến đẳng cấu chính
tắc.

Đặc biệt, từ hai bổ đề này suy ra rằng với mọi mở rộng đại số $K$ của $k$, và
mọi điểm hình học $a'$ của $X'=X\otimes_k K$ chiếu xuống điểm hình học $a$ của
$X$, đồng cấu có tính hàm tử
$\pi_1(X',a')\longrightarrow\pi_1(X,a)$ là đơn ánh.

\result{BỔ ĐỀ}{3} Cho $X$ là một lược đồ đầy đủ trên một vành Artin địa phương
$A$ sao cho $H^0(X,\Osh_X)=A$; cho $X'$ là một phủ không phân nhánh của $X$,
và đặt $A'=H^0(X',\Osh_{X'})$, vì thế là một vành hữu hạn trên $A$
(\textit{a priori} có thể phân nhánh trên $A$). Gọi $X_0,X'_0$ là các lược đồ
rút gọn liên kết với $X,X'$ (thu được bằng cách chia cho các bó phần tử lũy
linh của $\Osh_X,\Osh_{X'}$), và cho $k$ là một trường con của
$A/\mathfrak r(A)$ sao cho $A/\mathfrak r(A)$ hữu hạn trên trường con ấy (như vậy
$X_0$ là một lược đồ đại số đầy đủ trên $k$, còn $X'_0$ là một phủ không phân
nhánh của nó). Cuối cùng, cho $\Omega$ là một mở rộng đóng đại số của $k$, và
xét phủ không phân nhánh $X'_0\otimes_k\Omega$ của
$X_0\otimes_k\Omega$. Khi đó hai điều kiện sau tương đương:
\begin{enumerate}
  \item[\textit{i.}] $X'_0\otimes_k\Omega$ hoàn toàn phân rã trên
    $X_0\otimes_k\Omega$.
  \item[\textit{ii.}] Cấu xạ tự nhiên
    $X'\longrightarrow X\otimes_A A'$ là một đẳng cấu.
\end{enumerate}
Dưới các điều kiện ấy, $A'$ là một mở rộng không phân nhánh của $A$. Cuối
cùng, nếu $X'$ liên thông, điều kiện \textit{(i)} tương đương với điều kiện
sau đây, bề ngoài yếu hơn:
\begin{enumerate}
  \item[\textit{i bis.}] $X'_0\otimes_k\Omega$ có một tiết diện chính quy
    trên $X_0\otimes_k\Omega$.
\end{enumerate}
Khi điều kiện \textit{(ii)} được thỏa mãn, ta nói phủ không phân nhánh $X'$
của $X$ là \emph{tầm thường về mặt hình học}.

\src{212}
\result{BỔ ĐỀ}{4} Cho $f:X\longrightarrow Y$ là một cấu xạ proper sao cho
$f_*(\Osh_X)=\Osh_Y$. Cho $a$ là một điểm hình học của $X$ và $b$ là ảnh của
nó trên $Y$. Khi đó $\pi_1(X,a)\longrightarrow\pi_1(Y,b)$ là toàn ánh.

Thật vậy, cần chứng minh điều sau: nếu một phủ không phân nhánh $Y'$ của $Y$
(tương ứng với một bó đại số địa phương tự do $\mathcal A$) có tính chất
$X\otimes_Y Y'$ không liên thông, thì $Y'$ cũng không liên thông. Quả vậy,
$f^*(\mathcal A)$ khi ấy là tổng trực tiếp của hai bó vành khác không, nên ảnh
trực tiếp của nó cũng vậy; nhưng ảnh trực tiếp ấy chính là
$\mathcal A\otimes f_*(\Osh_X)=\mathcal A$.

\result{BỔ ĐỀ}{5} Cho $X$ là một lược đồ đầy đủ trên trường $k$; giả sử
$H^0(X,\Osh_X)$ là một vành địa phương $A$ và $A/\mathfrak r(A)$ là thuần bất
khả phân trên $k$. Cho $\Omega$ là một bao đóng đại số của $k$, và đặt
$\overline X=X\otimes_k\Omega$ (lược đồ này liên thông). Chọn một điểm hình
học $\overline a$ của $\overline X$ chiếu xuống điểm hình học $a$ của $X$;
khi đó có một dãy chính xác
\[
  e\longrightarrow\pi_1(\overline X,\overline a)
  \longrightarrow\pi_1(X,a)
  \longrightarrow\pi_1(k,b)\longrightarrow e
\]
(trong đó $\pi_1(k,b)$ là nhóm Galois của $\Omega$ trên $k$).

Tính đơn ánh của đồng cấu thứ nhất đã được chứng minh trong các Bổ đề 1 và 2;
tính chính xác ở giữa suy ra từ Bổ đề 3; cuối cùng, tính toàn ánh của đồng cấu
sau cùng (đây là chỗ duy nhất dùng việc $A/\mathfrak r(A)$ thuần bất khả phân)
suy ra từ Bổ đề 4.

\result{MỆNH ĐỀ}{4} Cho $f:X\longrightarrow Y$ là một cấu xạ proper và phẳng
sao cho, với mọi $y\in Y$, $H^0(f^{-1}(y),\Osh_{f^{-1}(y)})$ là một đại số khả
phân trên trường thặng dư $\kappa(y)$ (chẳng hạn, điều này đúng nếu
$f^{-1}(y)$ là một lược đồ khả phân trên $\kappa(y)$, tức là rút gọn và các
trường tương ứng với các thành phần bất khả quy của nó là những mở rộng khả
phân của $\kappa(y)$). Khi đó phủ $Y'$ của $Y$ liên kết với
$f_*(\Osh_X)$ là không phân nhánh.

Chứng minh suy ra dễ dàng từ Định lý 2.

\src{213}
Mệnh đề này, cùng với Bổ đề 1, trên thực tế quy việc nghiên cứu đồng luân của
các cấu xạ proper và phẳng (có thớ khả phân) về trường hợp
$f_*(\Osh_X)=\Osh_Y$ (vì khi dùng phân tích Stein, ta thay $Y$ bởi $Y'$).

\result{NHẬN XÉT}{} Một cấu xạ phẳng kiểu hữu hạn có các thớ là những lược đồ
khả phân (tương ứng, trơn) được gọi là \emph{khả phân} (tương ứng,
\emph{trơn}). Ta chứng minh được rằng nếu $f$ phẳng và $f^{-1}(y)$ khả phân
(tương ứng, trơn), thì tồn tại một lân cận của $f^{-1}(y)$ trên đó $f$ khả
phân (tương ứng, trơn). Kết quả tương tự đúng cho “chuẩn tắc tuyệt đối”
(“định lý Bertini”).

Cho $f:X\longrightarrow Y$ là một cấu xạ proper sao cho
\begin{equation}
  \tag{i}
  f_*(\Osh_X)=\Osh_Y,
\end{equation}
và cho $X'$ là một lược đồ hữu hạn trên $X$. Gọi $Y'$ là phủ của $Y$ tương ứng
với phân tích Stein của $X'\longrightarrow Y$ (xem Định lý 5). Cho $y\in Y$;
khi đó tập các thành phần liên thông của thớ $F'$ của $X'$ trên $y$ được đồng
nhất với tập các điểm $y'$ của $Y'$ nằm trên $y$ (Định lý 5). Xét cấu xạ hiển
nhiên
\begin{equation}
  \tag{*}
  X'\longrightarrow X\times_Y Y'
\end{equation}
suy ra từ các cấu xạ tự nhiên $X'\longrightarrow X$ và
$X'\longrightarrow Y'$\footnote{Bản in bỏ dấu phẩy trên $X'$ thứ hai; cấu xạ
đến $Y'$ là cấu xạ trong phân tích Stein của $X'\to Y$.}; cấu xạ này là một
đẳng cấu mỗi khi $X'$ có dạng $X\times_Y Y''$, trong đó $Y''$ là một phủ không
phân nhánh của $Y$; khi ấy $Y'$ chính là $Y''$ và (*) là đồng nhất. Ta muốn
đưa ra chính xác các điều kiện để $X'$ có dạng vừa nêu, tức là để $Y'$ không
phân nhánh và (*) là một đẳng cấu. Vì thế, xét thớ $F$ của $X$ trên $y$; đó là
một lược đồ proper trên $\kappa(y)$, còn $F'$ là một phủ của nó (không phân
nhánh nếu $X'$ không phân nhánh). Cho $F'_1$ là một thành phần liên thông của
$F'$ tương ứng với một điểm $y'_1$ của $Y'$ nằm trên $y$. Giả sử thêm rằng
\begin{enumerate}
  \item[\textit{(ii)}] $X'$ không phân nhánh trên $X$ tại các điểm của $F'_1$
    (do đó $F'_1$ là một phủ không phân nhánh của $F$), và
  \item[\textit{(iii)}] $F'_1$ là một phủ tầm thường về mặt hình học của $F$
    (xem Bổ đề 3).
\end{enumerate}

\result{ĐỊNH LÝ}{11} Dưới các điều kiện này, tồn tại một lân cận mở $U'$ của
$y'_1$ trong $Y'$ sao cho (*) là một đẳng cấu trên $U'$. Hơn nữa, $Y'$ không
phân nhánh tại $y'_1$ trên $Y$ (nhưng có thể phân nhánh tại các điểm $y'$ khác
của $Y'$ nằm trên $y$).

Dĩ nhiên, các điều kiện \textit{(ii)} và \textit{(iii)} cũng là cần thiết để
kết luận của định lý đúng. Chứng minh của định lý là dễ dàng, nhờ Bổ đề 3 và
Định lý 2.


\src{214}
\result{HỆ QUẢ}{1} Tiếp tục giả sử \textit{(i)} được thỏa mãn. Để một phủ
không phân nhánh của $X$ đẳng cấu với ảnh ngược của một phủ không phân
nhánh $Y'$ của $Y$, điều kiện cần và đủ là $X'$ cảm sinh trên mỗi thớ
$f^{-1}(y)$ một phủ tầm thường về mặt hình học.

Theo Định lý 11, tập các điểm của $Y$ tại đó điều kiện này được thỏa mãn là
mở, nên chỉ cần kiểm tra nó tại các điểm đóng $y$\ldots{} Ta ghi lại phát biểu
sau, tương đương với Hệ quả 1: hạt nhân của đồng cấu (toàn ánh theo Bổ đề 4)
$\pi_1(X)\longrightarrow\pi_1(Y)$ là nhóm con chuẩn tắc đóng sinh bởi các ảnh
trong $\pi_1(X)$ của các $\pi_1(\overline{f^{-1}(y)})$, trong đó
\[
  \overline{f^{-1}(y)}
  =f^{-1}(y)\otimes_{\kappa(y)}\overline{\kappa(y)},
\]
và $\overline{\kappa(y)}$ ký hiệu một bao đóng đại số của $\kappa(y)$. Cần lưu
ý rằng vì không thể chọn cùng một điểm gốc cho tất cả các thớ, các đồng cấu
$\pi_1(\overline{f^{-1}(y)})\longrightarrow\pi_1(X)$ trong mọi trường hợp chỉ
được xác định (sau khi đã chọn một điểm gốc cho $X$, và do đó cho $Y$) đến hợp
thành với một tự đẳng cấu trong của $\pi_1(X)$.

\result{HỆ QUẢ}{2} Dưới các giả thiết tổng quát của Định lý 11, giả sử thêm
rằng $Y,X,X'$ nguyên, và gọi $K,L,L'$ là các trường hàm tương ứng. Khi đó tồn
tại một mở rộng trung gian khả phân $K'$ của $K$ trong $L'$, độc lập tuyến
tính với $L$, sao cho $L'=LK'$ (do đó $L'=L\otimes_K K'$).

(Ta áp dụng phần cuối của Bổ đề 3 cho thớ hình học của $X$.) Trường hợp áp dụng
đáng quan tâm nhất của Định lý 11 là khi $f$ là một cấu xạ khả phân. Khi ấy
$X'$ cũng khả phân trên $Y$; do đó, theo Mệnh đề 4, $Y'$ không phân nhánh trên
$Y$, và vế thứ hai $X\times_Y Y'$ trong (*) không phân nhánh trên $X$. Ta dễ
dàng suy ra:

\result{HỆ QUẢ}{3} Ngoài \textit{(i)}, giả sử $f$ khả phân. Cho $X'$ là một
phủ không phân nhánh liên thông của $X$; để
$X'$\footnote{Bản in ghi “$X$” ở đây; chủ thể của tiêu chuẩn là phủ $X'$ được
đưa vào ở câu trước.} là ảnh ngược của một phủ không phân nhánh $Y'$ của $Y$,
điều kiện cần và đủ là phủ $\overline F'$ cảm sinh trên một thớ hình học
$\overline F=f^{-1}(y)$ có một tiết diện chính quy.

Cần lưu ý rằng không phải giả sử $\overline F'$ tầm thường về mặt hình học
trên $\overline F$ (điều này sẽ đúng \textit{a posteriori}, mặc dù
\textit{a priori} điều kiện ấy mạnh hơn nhiều). Thật ra, Hệ quả 3 tương đương
với phát biểu sau:

\src{215}
\result{HỆ QUẢ}{4} Cho $f:X\longrightarrow Y$ là một cấu xạ proper và khả
phân sao cho $f_*(\Osh_X)=\Osh_Y$; cho $\overline F$ là thớ hình học của một
điểm $y\in Y$, và chọn một điểm hình học trong $\overline F$. Qua các cấu xạ
$\overline F\longrightarrow X\longrightarrow Y$, ta thu được các điểm hình
học trong $X,Y$, được lấy làm điểm gốc cho các nhóm cơ bản của
$\overline F,X,Y$. Dưới các điều kiện ấy, ta có dãy chính xác
\[
  \boxed{
    \pi_1(\overline F)\longrightarrow\pi_1(X)
    \longrightarrow\pi_1(Y)\longrightarrow e
  }.
\]

Từ đó dễ dàng suy ra hai phát biểu sau của Serre--Lang, ở đây không còn cần
bất kỳ giả thiết chuẩn tắc nào:

\result{HỆ QUẢ}{5} Cho $X,Y$ là các lược đồ liên thông trên một trường $k$;
giả sử lược đồ rút gọn $X_{\mathrm{réd}}$ khả phân trên $k$ (điều này tự động
đúng nếu $k$ hoàn hảo), và giả sử $X$ hoặc $Y$ proper trên
$k$.\footnote{Điều kiện được bổ sung trong bản đính chính năm 1962.} Chọn một
điểm hình học $a$ (tương ứng, $b$) trong $X$ (tương ứng, $Y$), từ đó thu được
một điểm hình học $c=(a,b)$ trong $X\times_kY$ và một cấu xạ tự nhiên
\[
  \pi_1(X\times_kY,c)\longrightarrow
  \pi_1(X,a)\times\pi_1(Y,b)
\]
(suy ra từ các cấu xạ có tính hàm tử từ $\pi_1(X\times Y,c)$ vào
$\pi_1(X,a)$ và $\pi_1(Y,b)$). Cấu xạ này là đơn ánh, và thậm chí là song ánh
nếu $k$ đóng đại số.

(Tính toàn ánh trong trường hợp sau gần như hiển nhiên.) Cùng với Serre--Lang,
ta suy ra:

\result{HỆ QUẢ}{6} Cho $X$ là một lược đồ đại số liên thông và proper trên
một trường đóng đại số $k$, và cho $K$ là một mở rộng đóng đại số của $k$;
khi đó các nhóm cơ bản của $X$ và $X\otimes_kK$ là như nhau, tức là mọi phủ
không phân nhánh của lược đồ sau đều thu được bằng mở rộng vô hướng từ một phủ
không phân nhánh của $X$, duy nhất đến đẳng cấu.

\result{NHẬN XÉT}{}

1) Dùng Mệnh đề 4, ta thấy giả thiết $f_*(\Osh_X)=\Osh_Y$ trong Hệ quả 4 không
cốt yếu. Trong trường hợp tổng quát, thay vì đặt nhóm đơn vị $e$ sau
$\pi_1(Y)$, ta phải tiếp tục bằng
\[
  \pi_0(\overline F)\longrightarrow\pi_0(X)
  \longrightarrow\pi_0(Y)\longrightarrow e.
\]
như trong tôpô đại số.

\src{216}
2) Nói chung, hiện nay chưa thể nói gì về hạt nhân của
$\pi_1(\overline F)\longrightarrow\pi_1(X)$, nơi lẽ ra phải có sự tham gia của
một $\pi_2(Y)$. Tuy nhiên, có vẻ ta phải chứng minh được rằng
$\pi_1(\overline F)\longrightarrow\pi_1(X)$ là đơn ánh nếu $Y$ là phổ của một
vành địa phương $A$, bằng cách dựa vào Định lý 12 dưới đây (định lý này khẳng
định điều đó khi $A$ đầy đủ).

Định lý 11 chỉ dùng các Định lý 1 và 2. Bây giờ ta sẽ áp dụng Định lý 3, dựa
trên bổ đề sơ cấp sau:

\result{BỔ ĐỀ}{6} Cho $X$ là một lược đồ, và $X_0$ là lược đồ rút gọn tương
ứng (tức là lược đồ trong đó các phần tử lũy linh đã bị triệt tiêu). Khi đó mọi
phủ không phân nhánh $X'_0$ của $X_0$ đều được cảm sinh bởi một phủ không phân
nhánh $X'$ của $X$, được xác định duy nhất đến đẳng cấu chính tắc.

Bổ đề này, có bản chất hoàn toàn địa phương, đóng vai trò ở đây tương tự như
Định lý 8 trong lý thuyết môđun. Kết hợp nó với định lý tồn tại (Định lý 3),
ta suy ra:

\result{ĐỊNH LÝ}{12} Cho $A$ là một vành Noether địa phương đầy đủ có trường
thặng dư $k$. Cho $X$ là một lược đồ proper trên $A$. Khi đó mọi phủ không
phân nhánh $X'_0$ của $X_0=X\otimes_A k$ đều được cảm sinh bởi một phủ không
phân nhánh $X'$ của $X$, duy nhất đến đẳng cấu.

Nói cách khác:

\result{HỆ QUẢ}{1} Chọn một điểm hình học trong $X_0$ làm điểm gốc cho các
nhóm cơ bản của $X_0$ và $X$. Khi đó đồng cấu chính tắc
$\pi_1(X_0)\longrightarrow\pi_1(X)$ là một đẳng cấu.

Bây giờ áp dụng Bổ đề 5 cho $X_0$ (để đơn giản, giả sử
$H^0(X_0,\Osh_{X_0})=k$), và lưu ý rằng vì $A$ đầy đủ, các mở rộng không phân
nhánh của nó tương ứng với các mở rộng không phân nhánh của trường thặng dư;
tức là $\pi_1(Y)=\pi_1(k)$ (trong đó $Y=\operatorname{Spec}(A)$). Ta thu được
dãy chính xác:
\[
  e\longrightarrow\pi_1(\overline X_0)
  \longrightarrow\pi_1(X)\longrightarrow\pi_1(Y)\longrightarrow e.
\]

\result{HỆ QUẢ}{2} Cho $f:X\longrightarrow Y$ là một cấu xạ proper và phẳng,
cho $y_1$ là một điểm của $Y$ và $y_0$ là một chuyên hóa của $y_1$, và xét các
thớ “hình học” tương ứng $\overline X_1$ và $\overline X_0$; giả sử
$\overline X_0$ khả phân và liên thông (điều này kéo theo
$\overline X_1$ thỏa các điều kiện tương tự). Khi đó có thể tìm được một đồng
cấu nhóm
$\pi_1(\overline X_1)\longrightarrow\pi_1(\overline X_0)$, được xác định đến
một tự đẳng cấu trong, và đồng cấu này là toàn ánh.


\src{217}
Có thể hy vọng đồng cấu này luôn song ánh; tiếc thay, nói chung điều đó không
đúng nếu $\kappa(y_0)$ có đặc số $>0$. Tuy nhiên, dưới đây ta sẽ thu được một
chặn cho hạt nhân của đồng cấu ấy (ít nhất khi $\overline X_0$ trơn), từ đó suy
ra rằng nếu $\kappa(y_0)$ có đặc số $0$ thì đồng cấu trên là song ánh (một kết
quả cũng có thể chứng minh bằng phương pháp siêu việt). Dù sao, ta đã có một
chặn cho $\pi_1$ của một thớ đặc biệt theo nhóm tương ứng của một thớ tổng quát.
Chẳng hạn, dùng sự kiện một đường cong đại số trong đặc số $p$ có thể được nâng
lên một đường cong trong đặc số $0$ (Định lý 9, Hệ quả 4), ta thu được bằng
phương pháp siêu việt:

\result{HỆ QUẢ}{3} Cho $X_0$ là lược đồ của một đường cong đầy đủ và trơn trên
một trường đóng đại số có đặc số tùy ý, và gọi $g$ là giống của $X_0$; khi đó
$\pi_1(X_0)$ có $2g$ phần tử sinh tôpô, chịu ràng buộc bởi hệ thức quen thuộc.

Từ đó, bằng kỹ thuật tiết diện siêu phẳng quen thuộc, suy ra:

\result{HỆ QUẢ}{4} Cho $X$ là một lược đồ xạ ảnh trơn trên một trường đóng đại
số có đặc số tùy ý; khi đó $\pi_1(X)$ có hữu hạn phần tử sinh tôpô.

Ta tìm hạt nhân của đồng cấu
$\pi_1(\overline X_1)\longrightarrow\pi_1(\overline X_0)$. Để làm việc này,
có thể giả sử $Y$ là phổ của một vành định giá rời rạc đầy đủ
$V=\Osh_y$ ($y=y_0$). Câu hỏi tương đương với câu hỏi sau: cho một phủ không
phân nhánh $\overline X'_1$ của $\overline X_1$ (nếu muốn, có thể giả sử nó là
Galois), xác định những điều kiện để nó đến từ một phủ không phân nhánh của
$X_0$. \textit{A priori}, phủ đã cho đến từ, bằng mở rộng vô hướng, một phủ
không phân nhánh $X'_1$ của $X_1\otimes_K K'$, trong đó $K'$ là một mở rộng
con hữu hạn của bao đóng đại số $\overline K$ của trường phân thức $K$ của
$V$; nếu $\overline X'_1$ là Galois với nhóm $G$, ta có thể chọn $X'_1$ là
Galois với nhóm $G$. Khi đó ta có: để
$X'_1\otimes_{K'}\overline K=\overline X'_1$ đến từ một phủ không phân nhánh
của $X_0$, điều kiện cần và đủ là có thể tìm được một mở rộng hữu hạn $K''$
của $K'$ trong $\overline K$ sao cho
$X''_1=X'_1\otimes_{K'}K''$ có dạng $X''\otimes_{V''}K''$, trong đó $V''$ là
bao đóng nguyên của $V$ trong $K''$, còn $X''$ là một phủ không phân nhánh của
$X\otimes_VV''$. Chẳng hạn, giả sử $X_0$ chuẩn tắc tuyệt đối; khi đó
$X\otimes_VV''$ chuẩn tắc (vì nó phẳng trên $V''$ và có thớ đặc biệt chuẩn
tắc), và trường hàm của nó đồng nhất với trường $K''(X_1)$ của
\[
  X_1\otimes_KK''
  =(X\otimes_VK)\otimes_KK''
  =X\otimes_VK''
  =(X\otimes_VV'')\otimes_{V''}K''.
\]
Cho $L''=K''(X'_1)$ là trường hàm của
$X'_1\otimes_{K'}K''$; đây là một mở rộng hữu hạn khả phân của $K''(X_1)$, và
điều kiện trên cũng có nghĩa là $L''$ là một mở rộng không phân nhánh của
trường hàm của $X\otimes_VV''$ (tức là phép chuẩn hóa của
$X\otimes_VV''$ trong $L''$ không phân nhánh trên $X\otimes_VV''$). Hơn nữa,
chỉ cần kiểm tra rằng $L''$ không phân nhánh tại các điểm của thớ đặc biệt của
$X\otimes_VV''$ (vì nó không phân nhánh trên thớ tổng quát
$X_1\otimes_KK''$). Nếu bây giờ $X_0$ trơn, “định lý thuần khiết” của
Nagata--Zariski thậm chí cho thấy chỉ cần kiểm tra rằng $L''$ không phân nhánh
trên vành địa phương $O''$ tại điểm tổng quát của thớ đặc biệt của
$X\otimes_VV''$; đây là một vành định giá rời rạc, là phép chuẩn hóa trong
$K''(X_1)$ của vành địa phương $O\subset K(X_1)$ tại điểm tổng quát của thớ
đặc biệt của $X$. Do đó ta đã thu được:

\src{218}
\result{HỆ QUẢ}{5} Dưới các giả thiết và với các ký hiệu trên, để phủ không
phân nhánh $X'_1\otimes_{K'}\overline K$ của
$\overline X_1=X_1\otimes_K\overline K$ đến từ một phủ không phân nhánh của
$\overline X_0$, điều kiện cần và đủ là tồn tại một mở rộng con hữu hạn $K''$
của $\overline K/K'$ sao cho $K''(X'_1)$ không phân nhánh trên vành định giá
rời rạc $O''\subset K''(X_1)$.

Bây giờ lưu ý rằng $O''$ là phép chuẩn hóa trong $K''(X_1)$ của vành định giá
rời rạc $O'\subset K'(X_1)$ (phép chuẩn hóa của $O$ trong $K'(X_1)$), và
$O'$ chứa bao đóng nguyên $V'$ của $V$ trong $K'$; một phần tử đều hóa $u$ của
$V'$ cũng là một phần tử đều hóa của $O'$. Giả sử bây giờ $X'_1$ là Galois với
nhóm Galois $G$, có cấp $n$ nguyên tố cùng nhau với đặc số $p$ của
$\kappa(y_0)$ (đồng thời là đặc số của trường thặng dư của $O'$). Khi đó
$K'(X'_1)$ phân nhánh thuần trên $O'$, và dễ dàng suy ra (“bổ đề Abhyankar”)
rằng sau khi ghép thêm một căn bậc $n$, ký hiệu $v$, của một phần tử đều hóa
của $O'$, mở rộng ấy trở thành không phân nhánh trên phép chuẩn hóa của $O'$
trong $K'(X_1)(v)$. Nhưng có thể lấy $v$ là một căn bậc $n$ của một phần tử
đều hóa của $V'$, điều này chứng minh rằng điều kiện của Hệ quả 5 được thỏa
mãn. (Khả năng dùng bổ đề Abhyankar và định lý thuần khiết là do Serre bán cho
tôi.) Để diễn đạt kết quả thu được, với mọi nhóm compact hoàn toàn không liên
thông $\pi$, ta đưa vào nhóm thương $\widetilde\pi$ của $\pi$ theo nhóm con
đóng sinh bởi các nhóm con Sylow $p$ của nó, tức là giới hạn xạ ảnh của các
nhóm thương rời rạc của $\pi$ có cấp nguyên tố cùng nhau với $p$. Với ký hiệu
này, ta thu được:

\result{ĐỊNH LÝ}{13} Cho $f:X\longrightarrow Y$ là một cấu xạ proper và
phẳng, cho $y_1$ là một điểm của $Y$ và $y_0$ là một chuyên hóa của $y_1$, và
giả sử $\overline X_0$ liên thông và trơn. Khi đó đồng cấu
$\widetilde\pi_1(\overline X_1)\longrightarrow
\widetilde\pi_1(\overline X_0)$ suy ra từ đồng cấu toàn ánh trong Định lý 12,
Hệ quả 2, là một đẳng cấu.

Nói cách khác:

\src{219}
\result{HỆ QUẢ}{1} Việc phân loại các phủ Galois không phân nhánh có nhóm
Galois với cấp nguyên tố cùng nhau với đặc số $p$ của $\kappa(y_0)$ là như
nhau đối với $\overline X_0$ và $\overline X_1$.

Đặc biệt, nếu $\kappa(y_0)$ có đặc số không, bằng phương pháp đại số ta thu
được rằng $\pi_1(\overline X_1)\longrightarrow\pi_1(\overline X_0)$ là song
ánh.

Cuối cùng, ta lưu ý rằng các kỹ thuật đã dùng còn cho kết quả sau, tổng quát
hơn Định lý 13:

\result{ĐỊNH LÝ}{14} Cho $f:X\longrightarrow Y$ là một cấu xạ proper và trơn,
cho $D$ là một lược đồ con đóng của $X$, trơn trên $Y$ và có đối chiều 1 tại
mọi điểm của nó. Với một thớ $Z=f^{-1}(z)$ của $f$, đặt $Z'=Z-Z\cap D$ và gọi
$\pi_1^t(Z')$ là thương của nhóm cơ bản $\pi_1(Z')$ phân loại các phủ không
phân nhánh của $Z'$ phân nhánh thuần trên $\overline{Z\cap D}$. Cho $y_0,y_1$
như trong Định lý 13. Khi đó có một đồng cấu toàn ánh (được xác định đến một tự
đẳng cấu trong)
\[
  \pi_1^t(\overline X'_1)\longrightarrow\pi_1^t(\overline X'_0),
\]
và đồng cấu tương ứng
\[
  \widetilde\pi_1^t(\overline X'_1)
  \longrightarrow\widetilde\pi_1^t(\overline X'_0)
\]
là một đẳng cấu.

Từ đó suy ra các biến thể tương ứng của các hệ quả của Định lý 13 và của Hệ
quả 4 của Định lý 12. Tương tự, dùng Định lý 9, Hệ quả 3, bằng phương pháp
siêu việt ta thu được:

\result{HỆ QUẢ}{} Cho $X_0$ là lược đồ của một đường cong đầy đủ và trơn trên
một trường đóng đại số có đặc số tùy ý, và cho
$S=(s_i)_{1\le i\le n}$ là một tập con hữu hạn gồm $n$ phần tử của $X_0$. Khi
đó $\pi_1^t(X_0-S)$ có $2g+n$ phần tử sinh tôpô
$x_i,y_i$ ($1\le i\le g$) và $\sigma_j$ ($1\le j\le n$), chịu ràng buộc bởi
hệ thức
\[
  \left(\prod_i x_i y_i x_i^{-1}y_i^{-1}\right)
  \sigma_1\cdots\sigma_n=1,
\]
trong đó các $\sigma_j$ là những phần tử sinh của các nhóm quán tính tương ứng
với các $s_j$. Với mọi nhóm hữu hạn $G$ có cấp nguyên tố cùng nhau với đặc số,
được sinh bởi các phần tử $x_i,y_i,\sigma_j$ thỏa hệ thức trên, tồn tại một
phủ Galois không phân nhánh của $X_0-S$, với nhóm $G$, sao cho các nhóm quán
tính tại các điểm $s_j$ được sinh bởi các $\sigma_j$.

Khi $X_0$ có giống 0 và $n=3$, ta có một lời giải của “bài toán ba điểm”, ít
nhất đối với các phủ Galois có cấp nguyên tố cùng nhau với đặc số. (Ở đây,
Định lý 9 thật ra không cần dùng, và hơn nữa dường như hệ quả trước có thể được
suy ra từ trường hợp riêng xét trong bài toán ba điểm.)

\src{220}
\result{NHẬN XÉT}{}

1) Một nghiên cứu đầy đủ hơn, chắc chắn phải dùng các phủ Galois tổng quát hóa
của $X,X_0,X_1$ (với nhóm Galois có thể là vô cùng bé), hẳn sẽ cho phép thu
lại hạt nhân trong Định lý 12, Hệ quả 2. Trái lại, việc nghiên cứu các phủ cho
phép phân nhánh không “thuần” dường như phải khó hơn nhiều.

2) Bổ đề 6, kết hợp với một kết quả của Grauert về phép đầy đủ hóa hình thức
của một lược đồ xạ ảnh trơn dọc theo một tiết diện siêu phẳng (hoặc với định
lý, hiện chưa được chứng minh, nêu trong Nhận xét 2 sau Định lý 11), cũng cho
phép chứng minh trong hình học đại số “trừu tượng” định lý Lefschetz cổ điển về
nhóm cơ bản.

\begin{center}
  \textbf{TÀI LIỆU THAM KHẢO}
\end{center}

\begin{description}
  \item[{[1]}] Dieudonné (J.) và Grothendieck (A.). --- \emph{Éléments de
    géométrie algébrique}, Publications mathématiques de l'Institut des Hautes
    Études Scientifiques (sắp xuất bản).
  \item[{[2]}] Grothendieck (Alexander). --- \emph{The cohomology theory of
    abstract algebraic varieties}. International Congress of Mathematicians
    [1958. Edinburgh] (sắp xuất bản).
  \item[{[3]}] Serre (Jean-Pierre). --- \emph{Faisceaux algébriques
    cohérents}, Annals of Math., Series 2, t. 61, 1955, p. 197--278.
  \item[{[4]}] Serre (Jean-Pierre). --- \emph{Géométrie algébrique et
    géométrie analytique}, Ann. Institut Fourier Grenoble, t. 6, 1955--1956,
    p. 1--42.
\end{description}

\par
